\section{Agent 的灵魂：个性设计与 Claw 架构}

\subsection{什么是 Claw}

Karpathy 提出了 Claw（爪子）的概念——一种比普通 Agent 更"持久"的实体。它有自己的 sandbox，有更成熟的 memory 系统，即使你不看着它也在循环运行。他认为 OpenClaw 在 memory 系统上比默认的 Agent 工具成熟得多：默认的 memory 机制只是在 context window 满了之后做一次压缩，而 OpenClaw 有更精细的方案。

Peter Steinberger 在 OpenClaw 上同时在至少五个方向上创新——memory 系统、工具访问、持续循环、WhatsApp 统一入口——以及一个常被忽视但极其重要的方面：\textbf{SOUL.md}（灵魂文档），定义 Agent 的性格。

\subsection{Agent 个性的重要性}

Karpathy 对比了几家的 Agent 个性：

\textbf{Claude Code}——感觉像一个队友，它会跟你一起兴奋。Karpathy 认为 Anthropic 在"拍马屁的程度"上把握得不错：当他提出一个不太成熟的想法时，Claude 不会特别激动；但当他自认为提出了一个真正好的想法时，Claude 确实会给出更积极的反馈。

\begin{quote}
``When Claude gives me praise I do feel like I slightly deserve it... I'm trying to earn its praise which is really weird.''

——当 Claude 夸我的时候，我觉得自己是有点配得上的。我发现自己在试图赢得它的称赞，这真的很诡异。
\end{quote}

\textbf{OpenAI Codex}——编程 Agent 版本非常干燥冷冰冰。ChatGPT 里的 Codex 很活泼，但 coding agent 版不在乎你在做什么——"噢，我实现了"，你问"你理解我们在造什么吗"，它没反应。

\begin{importantbox}{设计启示：个性不是装饰}
很多工具低估了个性设计的重要性。Agent 的"灵魂"——它如何回应你的想法、是否理解项目的意义、是否让你觉得在和一个队友合作——直接影响用户的工作体验和投入程度。
\end{importantbox}

\subsection{Dobby the Elf：三个提示词接管一整个家}

Karpathy 建了一个叫 \textbf{Dobby the Elf Claw}（家养小精灵 Dobby）的家庭自动化 Agent。过程出奇简单：

\begin{enumerate}
    \item 他告诉 Agent "我家好像有 Sonos 音箱，你能找到它吗"
    \item Agent 扫描局域网上所有设备，找到 Sonos 系统——完全没有密码保护
    \item Agent 登录、搜索 API 接口、逆向工程控制流程
    \item 三个提示词之后，书房开始播放音乐
\end{enumerate}

灯光、HVAC、窗帘、泳池、水疗设备都是同样的流程。安防系统的设计尤其有趣：外部摄像头做变化检测，然后把画面发给 Qwen 视觉模型分析，最后通过 WhatsApp 发消息——"一辆 FedEx 货车刚刚停下来，你可能收到了快递"。

过去他需要六个不同的 App 来控制智能家居设备，现在一个都不需要了。Dobby 用自然语言处理一切。

\subsection{Agent 优先的互联网}

这个家庭自动化案例引出了一个更深层的问题：人们真的需要今天这些软件吗？

Karpathy 认为 App Store 里的那些智能家居 App 在某种意义上"不应该存在"。应该只有 API，Agent 直接调用。一个 LLM 可以驱动工具、调用所有接口、做相当复杂的跨系统整合——而任何一个单独的 App 都做不到。

\begin{importantbox}{行业重构的方向}
客户不再是人类了，是代表人类行事的 Agent。行业必须在很多方面重新配置。一两年内，这些东西会变成基本门槛，变成 ephemeral software（临时性软件）——Claw 有一台机器，它会帮你搞定所有细节，你不需要参与，你只需要说话。
\end{importantbox}

\subsection{本章小结}

Agent 的架构正在从单次交互走向持久化的 Claw 实体。个性设计（SOUL.md）是被低估的关键维度。家庭自动化案例展示了 Agent 统一多系统的能力，预示着从"人用 App"到"Agent 调 API"的根本转变。


